\documentclass[12pt]{article}
\usepackage{amsmath, amssymb, geometry, setspace, graphicx}
\geometry{margin=1in}
\setstretch{1.3}
\setlength{\parindent}{0pt}



\begin{document}




    {\LARGE \textbf{Inflation, Cost of Living, and Monetary Policy}}\\[0.5em]
    {\large ECON 204 -- Contemporary Economic Issues}\\[0.5em]
 Winter 2026\\


\hrule
\vspace{1em}



\section{Inflation and the Cost of Living}

Inflation refers to a sustained increase in the general price level of goods and services in an economy over time. When inflation occurs, the purchasing power of money falls, meaning that a given amount of income can buy fewer goods and services than before. It is important to distinguish between a one-time increase in prices and ongoing inflation. A temporary price increase, such as a rise in gasoline prices following a supply disruption, does not necessarily constitute inflation unless it leads to continuing price increases across many sectors of the economy. In contrast, persistent and widespread price growth indicates an inflationary environment. Canadian macroeconomics textbooks emphasize that inflation is best understood as a macroeconomic phenomenon that reflects economy-wide conditions rather than isolated price changes (Hubbard et al., Canadian Edition).

The cost of living refers to the amount of income required to maintain a given standard of living. In practice, economists measure changes in the cost of living using price indices, the most important of which is the Consumer Price Index (CPI). The CPI tracks the cost of purchasing a fixed basket of goods and services that represents typical household spending. This basket includes items such as food, housing, transportation, clothing, and various services. The CPI is constructed by comparing the cost of this basket in the current period to its cost in a base year. When the CPI rises, it indicates that the average price level has increased and that the cost of living has risen for households.

The inflation rate is calculated as the percentage change in the CPI from one period to the next. Formally, if $CPI_t$ represents the consumer price index in year $t$ and $CPI_{t-1}$ represents the index in the previous year, the inflation rate is given by:

\[
\pi_t = \frac{CPI_t - CPI_{t-1}}{CPI_{t-1}} \times 100
\]

For example, if the CPI increases from 158 to 163 over a year, the inflation rate is approximately 3.16 percent. This means that, on average, prices faced by consumers have increased by just over three percent.

While the CPI provides a useful measure of inflation, it does not perfectly capture changes in the true cost of living. One limitation is substitution bias. When the price of one good rises, consumers often substitute toward cheaper alternatives. A fixed basket of goods does not fully account for this behavior and may therefore overstate the increase in the cost of living. Another limitation is that the CPI may not fully reflect improvements in product quality. For instance, modern smartphones offer far more functionality than earlier models, even if their prices have not fallen significantly. Finally, the introduction of new goods and services, such as streaming platforms or digital technologies, may improve consumer welfare in ways that are not immediately captured by the CPI. Canadian textbooks stress that although the CPI is imperfect, it remains the most practical and widely used measure of inflation (Ragan, Canadian Edition).

\section{Types and Causes of Inflation}

Inflation can arise from different economic forces, and understanding these sources helps policymakers design appropriate responses. One major source of inflation is demand-pull inflation, which occurs when aggregate demand grows faster than the economy’s productive capacity. When households, firms, and governments increase their spending, total demand for goods and services rises. If the economy is operating near full capacity, firms respond by raising prices rather than expanding output. Expansionary fiscal policy, such as tax cuts or higher government spending, can contribute to demand-pull inflation by stimulating consumption and investment.

Another source of inflation is cost-push inflation, which results from increases in production costs. Rising wages, higher energy prices, or disruptions to supply chains can shift the short-run aggregate supply curve leftward. When firms face higher costs, they raise prices to maintain profitability. In this case, inflation is accompanied by slower economic growth or even recession, creating a difficult policy trade-off. The oil price shocks of the 1970s and recent global supply chain disruptions provide clear examples of cost-push inflation.

A third source of inflation arises from expectations. If workers and firms expect higher inflation in the future, they may adjust their behavior accordingly. Workers demand higher wages to protect their purchasing power, and firms raise prices to cover higher labor costs. This wage-price spiral can cause inflation to persist even if the original shock has passed. For this reason, modern central banks place great importance on managing inflation expectations.

\section{Inflation in the Aggregate Demand–Aggregate Supply Framework}

The aggregate demand–aggregate supply (AD–AS) model provides a useful framework for analyzing inflation. Aggregate demand represents the total spending on goods and services in the economy, while aggregate supply represents the total output firms are willing to produce at different price levels.

In a demand-pull inflation scenario, an increase in aggregate demand shifts the AD curve to the right. In the short run, this leads to both higher output and a higher price level. Over time, as wages and input costs rise, the short-run aggregate supply curve shifts leftward, bringing output back to its potential level but leaving the price level permanently higher. Inflation therefore results from sustained increases in aggregate demand relative to productive capacity.

In contrast, cost-push inflation occurs when the short-run aggregate supply curve shifts left due to higher production costs. This leads to higher prices and lower output simultaneously. Policymakers face a difficult choice in this situation because tightening monetary policy to reduce inflation can further weaken economic activity.

\section{Economic Consequences of Inflation}

Moderate and predictable inflation is not necessarily harmful. However, high or volatile inflation creates economic problems. One major issue is uncertainty. When firms and households cannot predict future prices, it becomes harder to plan investment, saving, and long-term contracts. Inflation also redistributes income and wealth in arbitrary ways. Borrowers benefit from unexpected inflation because they repay loans with money that is worth less in real terms, while lenders lose purchasing power. Households on fixed incomes, such as retirees without indexed pensions, may see their standard of living decline.

Inflation can also distort relative prices, making it harder for consumers and firms to distinguish between changes in overall prices and changes in the true scarcity of specific goods. These distortions reduce economic efficiency and can slow long-term growth.

\section{Nominal and Real Variables}

To understand the effects of inflation, economists distinguish between nominal and real variables. Nominal values are measured in current dollars, while real values are adjusted for inflation. For example, a nominal wage increase of four percent does not necessarily mean that workers are better off if inflation is three percent. In that case, real wages have only increased by about one percent.

The real interest rate is particularly important for monetary policy. It measures the true cost of borrowing in terms of purchasing power and is approximately equal to the nominal interest rate minus expected inflation:

\[
r \approx i - \pi^e
\]

If the nominal interest rate on a mortgage is six percent and expected inflation is two and a half percent, the real interest rate is about three and a half percent. This real rate influences decisions about consumption, saving, and investment.

\section{Monetary Policy in Canada}

In Canada, monetary policy is conducted by the Bank of Canada. The central objective of the Bank is to maintain price stability, which it interprets as keeping inflation low, stable, and predictable. Canada operates under an inflation-targeting framework, with a target of two percent inflation and a control range of one to three percent. This clear target helps anchor inflation expectations and enhances the credibility of monetary policy.

The Bank of Canada’s main policy instrument is the overnight interest rate. By raising or lowering this rate, the Bank influences borrowing costs throughout the financial system. Commercial banks adjust their lending rates, mortgage rates, and savings rates in response to changes in the policy rate.
\newpage
\section{The Monetary Transmission Mechanism}

Changes in the policy interest rate affect the economy through several channels. Higher interest rates discourage household consumption by making credit more expensive and encouraging saving. Business investment declines because borrowing becomes more costly and fewer projects remain profitable. Housing demand also falls as mortgage payments increase, slowing house price growth. In addition, higher Canadian interest rates can lead to an appreciation of the Canadian dollar, making imports cheaper and exports less competitive. This exchange rate channel further reduces inflationary pressures.

Through these channels, tighter monetary policy reduces aggregate demand and slows inflation. However, it may also reduce economic growth in the short run, illustrating the trade-offs faced by central banks.

\section{Inflation Expectations and Policy Credibility}

Inflation expectations play a central role in modern monetary policy. If households and firms believe that the Bank of Canada will successfully maintain inflation near its target, wage and price setting behavior will remain stable. This makes inflation easier to control and reduces the need for aggressive interest rate changes. Credibility therefore enhances the effectiveness of monetary policy and stabilizes the economy.

\section{Conclusion}

Inflation, the cost of living, and monetary policy are deeply interconnected in the Canadian economy. Inflation affects household purchasing power, business decisions, and income distribution. The CPI provides a practical measure of changes in the cost of living, even though it is imperfect. The Bank of Canada uses interest rates to influence economic activity and maintain price stability. By understanding the causes and consequences of inflation, students gain valuable insight into one of the most important challenges facing modern economies.
\newpage

\section*{References }

Hubbard, R. G., O’Brien, A. P., Serletis, A., et al. \textit{Macroeconomics}, Canadian Edition. Pearson.  

Ragan, C. \textit{The Economics of Macroeconomics}, Canadian Edition. Pearson.  

Parkin, M. \textit{Macroeconomics}, Canadian Edition. Pearson.  

Mankiw, N. G., Kneebone, R., and McKenzie, K. \textit{Principles of Macroeconomics}, Canadian Edition. Nelson.  

Mishkin, F. S., Serletis, A., and others. \textit{The Economics of Money, Banking, and Financial Markets}, 8th Canadian Edition. Pearson.


\end{document}
